\chapter{T-Duality, D-Branes, and Gauge Theory}

\lectureoverview{T-duality begins as an exact equivalence between closed
strings on circles of radii (R) and \(\alpha'/R\).  Momentum and winding
exchange roles, as do the worldsheet derivatives \(\partial_\tau Y\) and
\(\partial_\sigma Y\), and the spacetime vectors descending from
\(G_{\mu 5}\) and \(B_{\mu 5}\).  When the same operation is applied to an
open string, a Neumann endpoint condition becomes a Dirichlet condition.
The endpoint must then lie on a dynamical D-brane.  Open strings confined to
one brane produce an abelian gauge theory at low energy; strings joining a
stack of branes carry two endpoint labels and produce a nonabelian gauge
theory.  A D-string ending on a D3-brane supplies the corresponding magnetic
monopole.}

\section{Why Return to T-Duality?}

We have already met T-duality, but we return to it for a reason.  It was one
of the clues that string theory contained much more than perturbative
strings.  Once its consequences for open strings were understood, it led to
D-branes; D-branes in turn became a way to construct and study quantum field
theories.  The aim of this lecture is therefore not merely to repeat the
slogan that a small circle is equivalent to a large circle.  We shall rebuild
the argument carefully enough to see how that slogan forces new physical
objects into the theory.

The simplest setting is toroidal compactification.  A one-dimensional torus
is a circle.  A two-dimensional torus can be represented by a parallelogram
whose opposite sides are identified, and a three-dimensional torus by a
parallelepiped whose opposite faces are identified.  The same construction
continues in higher dimensions.  We shall keep only one compact spatial
coordinate, called \(Y\), and one representative noncompact coordinate,
called \(X\).  Time is present but is not shown in the board diagrams.

Choose the standard convention
\begin{equation}
  Y\sim Y+2\pi R,
  \label{eq:l10-circle-identification}
\end{equation}
so that \(R\) is the radius and \(2\pi R\) is the circumference.  The lecture
often calls the compact size simply \(r\) and suppresses factors of
\(2\pi\) and \(\alpha'\).  We shall retain the conventional factors in the
derivations and display the lecture units when they are useful.

For the first half of the discussion the strings are closed and oriented.
An orientation is an intrinsic arrow drawn around the string.  Reversing all
arrows everywhere would change no experiment; what matters is the relative
orientation of two strings.  During splitting and joining, the arrows must
join continuously.  This is the first hint that a signed, conserved quantity
will appear.

\begin{figure}[H]
  \centering
  \includegraphics[width=0.72\textwidth]{lecture_10_figure_01.png}
  \caption{Two closed strings with their relative orientations marked on
  the board.  Orientation must be continuous through splitting and joining.}
  \lectureframe{00:05:40}
  \label{fig:l10-oriented-closed-strings}
\end{figure}

\begin{classroomqa}
\audiencequestion{If the worldsheet coordinate runs from
\(\sigma=0\) to \(2\pi\), has the string already gone once around the
compact direction? \lecturetimestamp{00:08:42}}
\lecturerresponse{No.  The coordinate \(\sigma\) labels points along the
string; it does not label positions in the compact space.  Every closed
string has a periodic parameter \(\sigma\), whether it is wound or unwound.
The rubber-band analogy is useful: one may number the molecules around a
rubber band even when the band is not wrapped around anything.  Winding is a
property of the map \(Y(\tau,\sigma)\), not of the parameter \(\sigma\)
alone.}
\end{classroomqa}

\begin{classroomqa}
\audiencequestion{Must the compact size also be the same at every value of
the large coordinate \(X\)? \lecturetimestamp{00:10:00}}
\lecturerresponse{Not in a general dynamical compactification.  We begin
with a fixed circle because it is the cleanest background.  Later the
component \(G_{55}\) will become a scalar field whose value measures the
local size of the fifth direction, so that size can vary over ordinary
spacetime.}
\end{classroomqa}

\section{Winding Is Conserved Even When Strings Reconnect}

An unwound string may move through the compact direction without wrapping
it.  A wound string is different: as \(\sigma\) goes once around the closed
string, the embedding \(Y(\tau,\sigma)\) advances through one or more periods
of the compact circle.  Because the string is oriented, the advance has a
sign.  A string can wind in the positive direction or in the negative
direction.

The conservation law is clearest in the joining pictures.  Take one string
with winding \(+1\) and another with winding \(-1\).  Their compatible
segments may reconnect.  The resulting configuration has zero net winding
and can be deformed into an unwound string.  The board traces precisely this
opposite-orientation reconnection.

\begin{figure}[H]
  \centering
  \includegraphics[width=0.72\textwidth]{lecture_10_figure_02.png}
  \medskip

  \begin{tikzpicture}[x=1cm,y=1cm,>=Latex]
    \draw[thick] (-4.0,0.3) circle (0.65);
    \draw[->,thick] (-4.55,0.58) arc (150:-145:0.65);
    \node at (-4.0,-0.65) {$w=+1$};
    \draw[thick] (-2.0,0.3) circle (0.65);
    \draw[<-,thick] (-2.55,0.58) arc (150:-145:0.65);
    \node at (-2.0,-0.65) {$w=-1$};
    \node at (-3.0,0.3) {$+$};
    \draw[->,thick] (-1.0,0.3) -- (0.3,0.3);
    \draw[thick] (1.4,0.3) ellipse (0.95 and 0.55);
    \node at (1.4,-0.65) {$w_{\rm total}=0$};
  \end{tikzpicture}
  \caption{Opposite windings reconnect to a zero-net-winding
  configuration.  The lower diagram is a clean reconstruction of the board
  move, not a second dynamical process.}
  \lectureframe{00:13:20}
  \label{fig:l10-opposite-winding}
\end{figure}

Now repeat the experiment with two strings having the same orientation.
They can join, but they cannot unwind.  The result carries twice the winding
of either incoming string.  The number of connected strings has changed;
the total winding has not.

\begin{figure}[H]
  \centering
  \includegraphics[width=0.72\textwidth]{lecture_10_figure_03.png}
  \medskip

  \begin{tikzpicture}[x=1cm,y=1cm,>=Latex]
    \draw[thick] (-4.0,0.3) circle (0.65);
    \draw[->,thick] (-4.55,0.58) arc (150:-145:0.65);
    \node at (-4.0,-0.65) {$w=+1$};
    \draw[thick] (-2.0,0.3) circle (0.65);
    \draw[->,thick] (-2.55,0.58) arc (150:-145:0.65);
    \node at (-2.0,-0.65) {$w=+1$};
    \node at (-3.0,0.3) {$+$};
    \draw[->,thick] (-1.0,0.3) -- (0.3,0.3);
    \draw[thick] (1.45,0.3) circle (0.78);
    \draw[thick] (1.45,0.3) circle (0.52);
    \draw[->,thick] (0.81,0.74) arc (145:-145:0.78);
    \node at (1.45,-0.78) {$w=+2$};
  \end{tikzpicture}
  \caption{Same-sign windings can reconnect into a doubly wound string.
  Connectivity changes while total winding remains two.}
  \lectureframe{00:16:40}
  \label{fig:l10-same-winding}
\end{figure}

This distinction survives every allowed process.  A total winding of one
hundred might be carried by one string wound one hundred times, two strings
wound fifty times each, or one hundred singly wound strings.  Strings can
split and join among these descriptions, but the signed sum remains one
hundred.  We denote it by \(w\in\mathbb Z\).

\section{Momentum Modes and Winding Modes}

Conservation by itself is not yet duality.  We next compare energies.  A
state moving around the compact circle has quantized momentum
\begin{equation}
  p_Y=\frac{n}{R},
  \qquad n\in\mathbb Z.
  \label{eq:l10-compact-momentum}
\end{equation}
For a state whose other mass contributions can be neglected, the compact
motion contributes \(E_{\rm mom}=|n|/R\).  Positive and negative \(n\)
denote opposite directions of motion, while the energy is positive.

\begin{figure}[H]
  \centering
  \includegraphics[width=0.72\textwidth]{lecture_10_figure_04.png}
  \caption{The compact momentum quantum number on the board.  The lecture
  writes the momentum schematically as an integral of the compact velocity.}
  \lectureframe{00:21:30}
  \label{fig:l10-compact-momentum-board}
\end{figure}

A string wound \(|w|\) times has length \(2\pi |w|R\).  With fundamental
string tension
\begin{equation}
  T_{\rm F}=\frac{1}{2\pi\alpha'},
\end{equation}
its winding energy is
\begin{equation}
  E_{\rm wind}=T_{\rm F}(2\pi |w|R)
  =\frac{|w|R}{\alpha'}.
  \label{eq:l10-winding-energy}
\end{equation}
In the lecture units \(\alpha'=1\), this is simply \(E_{\rm wind}=|w|R\).

The two towers behave oppositely as the radius changes.  For \(R\gg
\sqrt{\alpha'}\), momentum levels are closely spaced and winding levels are
widely spaced.  For \(R\ll\sqrt{\alpha'}\), momentum levels are widely
spaced and winding levels are closely spaced.

\begin{figure}[H]
  \centering
  \begin{minipage}[t]{0.48\textwidth}
    \centering
    \includegraphics[width=\linewidth]{lecture_10_figure_05.png}
    \par\smallskip
    \textit{Large compact radius.}
  \end{minipage}\hfill
  \begin{minipage}[t]{0.48\textwidth}
    \centering
    \includegraphics[width=\linewidth]{lecture_10_figure_06.png}
    \par\smallskip
    \textit{Small compact radius.}
  \end{minipage}
  \caption{The two spectral comparisons on the board.  Sparse and dense
  towers exchange roles as the radius is inverted.}
  \lectureframe{00:25:30--00:26:00}
  \label{fig:l10-dual-spectra-board}
\end{figure}

The clean schematic is
\begin{figure}[H]
  \centering
  \begin{adjustbox}{max width=\linewidth}
  \begin{tikzpicture}[x=1cm,y=0.74cm,>=Latex]
    \node[font=\bfseries] at (-3.6,6.2) {$R\gg\sqrt{\alpha'}$};
    \draw[->] (-5.1,0) -- (-5.1,5.5) node[above] {$E$};
    \foreach \y in {0.4,1.0,1.6,2.2,2.8,3.4,4.0,4.6}
      \draw (-5.35,\y)--(-4.85,\y);
    \node[align=center] at (-4.25,5.15) {momentum\\dense};
    \draw[->] (-2.2,0) -- (-2.2,5.5) node[above] {$E$};
    \foreach \y in {0.4,2.4,4.4}
      \draw (-2.45,\y)--(-1.95,\y);
    \node[align=center] at (-1.35,5.15) {winding\\sparse};

    \node[font=\bfseries] at (3.6,6.2) {$R\ll\sqrt{\alpha'}$};
    \draw[->] (2.2,0) -- (2.2,5.5) node[above] {$E$};
    \foreach \y in {0.4,2.4,4.4}
      \draw (1.95,\y)--(2.45,\y);
    \node[align=center] at (3.05,5.15) {momentum\\sparse};
    \draw[->] (5.1,0) -- (5.1,5.5) node[above] {$E$};
    \foreach \y in {0.4,1.0,1.6,2.2,2.8,3.4,4.0,4.6}
      \draw (4.85,\y)--(5.35,\y);
    \node[align=center] at (5.95,5.15) {winding\\dense};
  \end{tikzpicture}
  \end{adjustbox}
  \caption{Reconstruction of the level-spacing argument.  The complete
  string spectrum is unchanged when momentum and winding are exchanged
  together with \(R\leftrightarrow\alpha'/R\).}
  \label{fig:l10-dual-spectra-redraw}
\end{figure}

Suppose an observer is handed one of these spectra without being told how it
was made.  Is the closely spaced tower momentum on a large circle, or
winding on a small circle?  The spectrum alone cannot say.  String theory
makes the stronger statement that scattering amplitudes cannot say either.
The two descriptions are exactly equivalent under
\begin{equation}
  n\longleftrightarrow w,
  \qquad
  R\longleftrightarrow \widetilde R=\frac{\alpha'}{R}.
  \label{eq:l10-radius-duality}
\end{equation}
The fixed point is
\begin{equation}
  R=\sqrt{\alpha'},
  \label{eq:l10-self-dual-radius}
\end{equation}
or \(R=1\) in the lecture units.  At that radius additional states can
become massless and the gauge symmetry is enhanced.

\begin{classroomqa}
\audiencequestion{Is winding only a mathematical convenience for rewriting
the spectrum? \lecturetimestamp{00:27:33}}
\lecturerresponse{No.  Winding labels real sectors of the closed-string
Hilbert space, and those sectors take part in interactions.  The word
``duality'' means that the momentum description and winding description are
different coordinates on the same physics, not that one sector is an
optional mnemonic for the other.}
\end{classroomqa}

\begin{classroomqa}
\audiencequestion{Is there a radius that is mapped to itself?
\lecturetimestamp{00:27:58}}
\lecturerresponse{Yes.  The self-dual value is
\(R=\sqrt{\alpha'}\), which is \(R=1\) in the units used on the board.
The two towers meet there, and the spectrum has enhanced symmetry.}
\end{classroomqa}

\section{The Worldsheet Form of the Duality}

The spectral argument is suggestive, but the worldsheet makes the exchange
more precise.  In conformal gauge, the momentum conjugate to \(Y\) is
proportional to \(\partial_\tau Y\).  Thus the total compact momentum has
the form
\begin{equation}
  P_Y=\frac{1}{2\pi\alpha'}
      \int_0^{2\pi}d\sigma\,\partial_\tau Y.
  \label{eq:l10-worldsheet-momentum}
\end{equation}
The winding number is
\begin{equation}
  w=\frac{1}{2\pi R}
      \int_0^{2\pi}d\sigma\,\partial_\sigma Y
   =\frac{Y(2\pi)-Y(0)}{2\pi R}.
  \label{eq:l10-worldsheet-winding}
\end{equation}

\begin{classroomqa}
\audiencequestion{Why does the integral of a derivative not vanish around
the closed string? \lecturetimestamp{00:32:33}}
\lecturerresponse{Because \(Y\) is a periodic target-space coordinate, not
a single-valued real function on the \(\sigma\)-circle.  The worldsheet
point \(\sigma=2\pi\) is the same as \(\sigma=0\), but a lift of the
embedding to the covering line changes by \(2\pi wR\).  The integral of
\(\partial_\sigma Y\) measures that shift.}
\end{classroomqa}

\begin{figure}[H]
  \centering
  \includegraphics[width=0.76\textwidth]{lecture_10_figure_07.png}
  \caption{The momentum and winding integrals after both lines have been
  completed on the board.  Common \(2\pi\) and \(\alpha'\) factors are
  suppressed in the lecture normalization.}
  \lectureframe{00:33:25}
  \label{fig:l10-momentum-winding-integrals}
\end{figure}

For a once-wound string,
\begin{equation}
  Y(\tau,\sigma+2\pi)=Y(\tau,\sigma)+2\pi R.
\end{equation}
For winding \(w\), the right-hand shift is \(2\pi wR\).  Equation
\eqref{eq:l10-worldsheet-winding} then counts the winding exactly.  This is
the precise version of the board shorthand \(w\sim R^{-1}\int
d\sigma\,\partial_\sigma Y\).

The worldsheet statement of T-duality is therefore
\begin{equation}
  \partial_\tau Y\longleftrightarrow
  \partial_\sigma \widetilde Y,
  \qquad
  \partial_\sigma Y\longleftrightarrow
  \partial_\tau \widetilde Y,
  \label{eq:l10-derivative-duality}
\end{equation}
with momentum in \(Y\) becoming winding in the dual coordinate
\(\widetilde Y\), and conversely.

\begin{classroomqa}
\audiencequestion{Does this mean that the worldsheet coordinates
\(\tau\) and \(\sigma\) themselves are related or interchanged?
\lecturetimestamp{00:34:39}}
\lecturerresponse{No.  T-duality maps solutions and exchanges particular
derivatives of the compact target coordinate.  It does not identify the
worldsheet time coordinate with the worldsheet spatial coordinate.}
\end{classroomqa}

\section{Why Winding States Cannot Be Avoided}

One might ask why a string would ever be found in a wound state when an
unwound state is available.  A string does not choose its state by
preference.  Energy supplied in a collision stretches it.  A sufficiently
energetic, highly distorted string can reconnect into two strings with
opposite winding.  Their total winding remains zero, but after the split the
two strings may separate.  One member of the pair can be sent far away,
leaving an observer with an apparently isolated wound string.

\begin{figure}[H]
  \centering
  \includegraphics[width=0.76\textwidth]{lecture_10_figure_08.png}
  \caption{A highly stretched zero-net-winding configuration from which
  oppositely wound strings can separate after reconnection.}
  \lectureframe{00:36:15}
  \label{fig:l10-winding-pair-production}
\end{figure}

The energy threshold is set by the winding energy.  On a large circle it
takes a large collision energy to stretch enough string around the compact
direction.  Above threshold, however, winding--antiwinding pair production
cannot be forbidden.  The analogy is photon collisions producing an
electron--positron pair.  The total electric charge is still zero, yet the
charged particles are present and can be separated.  Winding conservation
works in the same way.

This analogy is more than verbal.  Compact momentum and winding behave as
two distinct charge quantum numbers.  Like signs repel and opposite signs
attract in the corresponding reduced description.  We should therefore ask
for the two vector fields that mediate those forces.

\section{Momentum Charge from Five-Dimensional Gravity}

Begin with gravity in five dimensions.  Let capital indices \(M,N\) run
over the four ordinary spacetime directions and the compact fifth
direction.  The metric decomposes as
\begin{equation}
  G_{MN}\quad\longrightarrow\quad
  \bigl\{G_{\mu\nu},\,G_{\mu5},\,G_{55}\bigr\}.
  \label{eq:l10-metric-decomposition}
\end{equation}
The symmetric component \(G_{5\mu}\) contains no new information beyond
\(G_{\mu5}\).

\begin{figure}[H]
  \centering
  \includegraphics[width=0.74\textwidth]{lecture_10_figure_09.png}
  \caption{The independent groups of five-dimensional metric components:
  \(G_{\mu\nu}\), \(G_{\mu5}\), and \(G_{55}\).}
  \lectureframe{00:44:05}
  \label{fig:l10-metric-components}
\end{figure}

From the four-dimensional point of view,
\begin{align}
  G_{\mu\nu} &\quad\hbox{is a spacetime metric},\\
  G_{\mu5} &\equiv A_\mu\quad\hbox{is a vector field},\\
  G_{55} &\equiv e^{2\varphi}\quad\hbox{is a scalar modulus}.
  \label{eq:l10-four-dimensional-fields}
\end{align}
The scalar \(\varphi\) measures the size of the compact direction.  It can
vary smoothly over ordinary spacetime, so waves of \(\varphi\) are waves in
the radius of the fifth dimension.

\begin{figure}[H]
  \centering
  \includegraphics[width=0.74\textwidth]{lecture_10_figure_10.png}
  \caption{The board identifies \(G_{\mu5}\) with a vector \(A_\mu\) and
  \(G_{55}\) with a scalar \(\Phi\).}
  \lectureframe{00:47:40}
  \label{fig:l10-vector-and-radius-modulus}
\end{figure}

The source of a mixed gravitational component is mixed momentum flow.  In
particular, momentum around the fifth direction sources \(G_{\mu5}\).  A
four-dimensional observer therefore interprets the integer \(n\) as an
electric charge coupled to \(A_\mu\).  This is the Kaluza--Klein mechanism:
what is momentum in the higher-dimensional geometry becomes charge in the
lower-dimensional field theory.

\editorialnote{The lecture loosely calls the scalar from \(G_{55}\) a
``dilaton.''  In this reduction it is specifically the radius modulus.  The
ten-dimensional string dilaton, which controls the string coupling, is a
separate field, although lower-dimensional scalar combinations can mix
under changes of frame and duality.}

\begin{classroomqa}
\audiencequestion{Is the scalar \(\varphi\) proportional to the momentum
integer \(n\)? \lecturetimestamp{00:48:43}}
\lecturerresponse{No.  The integer \(n\) labels a quantized state and acts
as charge.  The scalar \(\varphi(x)\) is a continuous field that determines
the local compactification size.  It may vary from place to place; it is not
itself an integer.}
\end{classroomqa}

\section{The Second Vector Field and Winding Charge}

Winding is also charge-like, but it cannot couple to the same vector in the
same way.  To see the second vector, return to the first closed-string
excitations.  A closed string has left-moving and right-moving oscillator
modes.  Physical states obey level matching: the total left-moving and
right-moving excitation levels must agree.

In the bosonic theory the unexcited state is the tachyon.  A single
left-moving or right-moving oscillator would violate level matching.  The
first allowed massless state therefore contains one oscillator of each
kind,
\begin{equation}
  \alpha_{-1}^{I}\widetilde\alpha_{-1}^{J}\ket{0;k}.
  \label{eq:l10-first-massless-state}
\end{equation}
The indices \(I,J\) run over transverse spatial directions, including the
compact direction.  If one index is ordinary and the other is \(5\), there
are two vector-like oscillator states,
\begin{equation}
  \alpha_{-1}^{\mu}\widetilde\alpha_{-1}^{5}\ket{0;k},
  \qquad
  \alpha_{-1}^{5}\widetilde\alpha_{-1}^{\mu}\ket{0;k}.
  \label{eq:l10-two-vector-states}
\end{equation}

\begin{figure}[H]
  \centering
  \includegraphics[width=0.76\textwidth]{lecture_10_figure_11.png}
  \caption{The two left--right oscillator orderings that become
  four-dimensional vector states when one polarization points along the
  compact direction.}
  \lectureframe{00:58:20}
  \label{fig:l10-two-vector-states-board}
\end{figure}

Their symmetric and antisymmetric combinations belong to different
spacetime fields:
\begin{align}
  \alpha_{-1}^{(\mu}\widetilde\alpha_{-1}^{5)}\ket{0;k}
    &\longleftrightarrow G_{\mu5},\\
  \alpha_{-1}^{[\mu}\widetilde\alpha_{-1}^{5]}\ket{0;k}
    &\longleftrightarrow B_{\mu5}.
  \label{eq:l10-g-and-b-vectors}
\end{align}
The first is the Kaluza--Klein vector sourced by compact momentum.  The
second descends from the antisymmetric Kalb--Ramond field and is sourced by
string winding.  Thus the reduced theory contains two photon-like fields,
one for each charge sector.

If both oscillator indices point along the compact direction, the state is
a scalar associated with the compactification modulus.  If both are in
ordinary transverse directions, the symmetric traceless part is the
graviton, the antisymmetric part is \(B_{\mu\nu}\), and the trace contributes
to a scalar sector.

\begin{classroomqa}
\audiencequestion{Which of these states is the singlet and which is the
triplet? \lecturetimestamp{01:00:21}}
\lecturerresponse{That classification is not the useful one here.  Both
mixed states transform as four-dimensional vectors.  Their sum and
difference separate the symmetric metric polarization from the
antisymmetric \(B\)-field polarization.  The state with two compact indices
is instead a four-dimensional scalar.}
\end{classroomqa}

The complete dictionary reached before the break is
\begin{equation}
  \boxed{
  \begin{gathered}
    n\leftrightarrow w,
    \qquad
    R\leftrightarrow\dfrac{\alpha'}{R},\\
    \partial_\tau Y\leftrightarrow\partial_\sigma\widetilde Y,
    \qquad
    G_{\mu5}\leftrightarrow B_{\mu5}.
  \end{gathered}}
  \label{eq:l10-complete-t-duality-map}
\end{equation}

\begin{figure}[H]
  \centering
  \includegraphics[width=0.72\textwidth]{lecture_10_figure_12.png}
  \caption{The board summary of the T-duality map: momentum and winding,
  radius and inverse radius, worldsheet derivatives, and the two vector
  fields are exchanged.  The lecturer partly obscures the final line, which
  is therefore supplied explicitly in Equation
  \eqref{eq:l10-complete-t-duality-map}.}
  \lectureframe{01:02:20}
  \label{fig:l10-t-duality-map-board}
\end{figure}

This construction should not be confused with a claim that every field just
listed has already been identified in the observed particle spectrum.  At
this point it is a mathematical consequence of the compactified string
theory.  There are many proposals for connecting such structures to
phenomenology, but that is not the problem being solved in this lecture.

\section{Open Strings Change the Question}

After the break, we begin again with open strings.  A closed string can have
stable winding because it has no endpoint through which the winding can
slip away.  An open string has endpoints and therefore no analogous
topologically protected winding number.  Yet T-duality is an exact symmetry
of the theory.  What does the symmetry do to an open string?

For a freely moving endpoint, the variation of the string action gives a
Neumann boundary condition.  In the simple \((X,Y)\) geometry,
\begin{equation}
  \left.\partial_\sigma X\right|_{\partial\Sigma}=0,
  \qquad
  \left.\partial_\sigma Y\right|_{\partial\Sigma}=0.
  \label{eq:l10-neumann-boundary}
\end{equation}
The mechanical picture is straightforward.  Regard the string as a chain
of ever smaller mass elements.  A nonzero tangent force on the last element
would give it an acceleration that diverges as its mass tends to zero.  A
free endpoint therefore adjusts until the net endpoint tension vanishes.
Someone could instead hold the endpoint fixed and supply the necessary
force, but then the endpoint would no longer be free.

\begin{figure}[H]
  \centering
  \includegraphics[width=0.76\textwidth]{lecture_10_figure_13.png}
  \caption{The open string and its completed Neumann endpoint conditions,
  \(\partial_\sigma X=\partial_\sigma Y=0\).}
  \lectureframe{01:09:20}
  \label{fig:l10-neumann-board}
\end{figure}

Now dualize only the compact coordinate \(Y\).  The noncompact coordinate
\(X\) is unchanged.  Equation \eqref{eq:l10-derivative-duality} turns the
second Neumann condition into
\begin{equation}
  \left.\partial_\tau\widetilde Y\right|_{\partial\Sigma}=0.
  \label{eq:l10-dirichlet-time-derivative}
\end{equation}
Therefore the endpoint has a time-independent position in the dual compact
direction,
\begin{equation}
  \left.\widetilde Y\right|_{\partial\Sigma}=Y_0.
  \label{eq:l10-dirichlet-boundary}
\end{equation}
This is a Dirichlet boundary condition.

\begin{figure}[H]
  \centering
  \includegraphics[width=0.76\textwidth]{lecture_10_figure_14.png}
  \medskip

  \begin{adjustbox}{max width=\linewidth}
  \begin{tikzpicture}[x=1cm,y=1cm,>=Latex]
    \node[align=center] at (-3.5,0.8)
      {Neumann\\\(\left.\partial_\sigma Y\right|_{\partial\Sigma}=0\)};
    \draw[->,thick] (-1.7,0.8) -- (1.3,0.8)
      node[midway,above] {$T_Y$};
    \node[align=center] at (3.4,0.8)
      {Dirichlet\\\(\left.\partial_\tau\widetilde Y\right|_{\partial\Sigma}=0\)};
    \node[align=center] at (3.4,-0.5)
      {\(\widetilde Y(\tau)=Y_0\) at each endpoint};
  \end{tikzpicture}
  \end{adjustbox}
  \caption{The completed dual endpoint condition on the board and its clean
  logical reconstruction.  T-duality turns freedom in the compact direction
  into a fixed endpoint position in the dual direction.}
  \lectureframe{01:12:35}
  \label{fig:l10-neumann-to-dirichlet}
\end{figure}

\begin{classroomqa}
\audiencequestion{Why does the dual condition mean that the endpoint is
fixed? \lecturetimestamp{01:12:39}}
\lecturerresponse{Because \(\partial_\tau\widetilde Y=0\) is zero velocity
in the dualized direction.  Wherever the endpoint begins in
\(\widetilde Y\), it remains there.  Its coordinates parallel to the
un-dualized directions are still free.}
\end{classroomqa}

\section{The Object Holding the Endpoint Is a D-Brane}

If an endpoint is fixed, something must be able to hold it there.  Exact
T-duality therefore requires a new object on which open strings can end.
That object is a D-brane.  The letter \(D\) stands for Dirichlet, not for the
dimension of the object.  In the historical aside of the lecture, the
objects might have been called after Joseph Polchinski, whose analysis made
their dynamical role unmistakable, but the already established term
``p-brane'' left Dirichlet as the natural name.

A \(p\)-brane has \(p\) spatial dimensions.  A string is a one-brane, a
membrane is a two-brane, and an extended three-dimensional object is a
three-brane.  A D\(p\)-brane is a \(p\)-brane on which fundamental-string
endpoints obey Dirichlet conditions in the transverse directions and
Neumann conditions in the directions tangent to the brane.

Nothing in the duality argument singles out one value of \(Y_0\), so the
D-brane cannot be a rigid fixture nailed to a preferred point in the
universe.  It must be movable.  Relativity also forbids an absolutely rigid
extended body, so a D-brane can fluctuate and bend.  The simple flat brane
drawn on the board is its equilibrium configuration, not a restriction on
its dynamics.

\begin{classroomqa}
\audiencequestion{Must the two endpoints of an open string have the same
fixed value of the dual coordinate? \lecturetimestamp{01:14:55}}
\lecturerresponse{Not in general.  Both endpoints may lie on the same
D-brane, in which case they share that brane's transverse location.  They
may also lie on two different D-branes at different transverse positions,
and the string then stretches between them.  A brane itself may fluctuate
and bend, so the endpoint locus need not remain a perfectly flat plane.}
\end{classroomqa}

\section{Counting the Dimensions of a D-Brane}

Ten-dimensional superstring theory has nine spatial dimensions.  Suppose
T-duality imposes a Dirichlet condition in one spatial direction.  The
endpoint is still free in the other eight, so its allowed locus is an
eight-dimensional D8-brane.  Each additional independent Dirichlet
condition removes one direction of endpoint motion.  If \(k\) spatial
coordinates are fixed, then
\begin{equation}
  p=9-k.
  \label{eq:l10-dp-counting}
\end{equation}

The tabletop gives the three-dimensional version of the same count.  Fix
the vertical coordinate and a fingertip may move on a two-dimensional
tabletop.  Fix a second coordinate and it may move only along a line.  Fix
all three and it is confined to a point.

\begin{figure}[H]
  \centering
  \begin{adjustbox}{max width=\linewidth}
  \begin{tikzpicture}[x=1.05cm,y=1cm,>=Latex]
    \draw[fill=blue!8] (-5.5,-0.7) -- (-2.5,-0.7) -- (-1.7,0.2)
      -- (-4.7,0.2) -- cycle;
    \node at (-3.6,-1.15) {one constraint: sheet};
    \draw[very thick,blue!60!black] (-0.9,-0.25) -- (1.3,-0.25);
    \node at (0.2,-1.15) {two constraints: line};
    \fill[blue!70!black] (4.1,-0.25) circle (0.09);
    \node at (4.1,-1.15) {three constraints: point};
    \draw[->] (-1.45,0.95) -- (-0.75,0.95);
    \draw[->] (2.25,0.95) -- (2.95,0.95);
    \node at (0.35,1.35) {each Dirichlet condition lowers \(p\) by one};
  \end{tikzpicture}
  \end{adjustbox}
  \caption{The dimensional-counting argument used at the tabletop.  A
  D\(p\)-brane is the locus in which endpoints remain free to move.}
  \label{fig:l10-dp-counting-redraw}
\end{figure}

Repeated dualities generate D8-, D7-, D6-branes and so on down to a D0-brane,
which is pointlike in space.  If no spatial direction is fixed, a D9-brane
fills all nine spatial dimensions and ordinary open strings are free to move
everywhere.  The lecture briefly mentions D\((-1)\)-branes; these are
localized in Euclidean spacetime and behave as instantons rather than as
ordinary particles propagating through time.

\editorialnote{The dimensional ladder is a statement about the type-II
theories taken together.  Stable BPS D\(p\)-branes have even \(p\) in type
IIA and odd \(p\) in type IIB.  T-duality along one circle changes \(p\) by
one and simultaneously exchanges type IIA with type IIB.}

\begin{classroomqa}
\audiencequestion{If only six dimensions are compact, should the counting
start by subtracting from six rather than nine? \lecturetimestamp{01:20:49}}
\lecturerresponse{The elementary derivation performs T-duality on compact
directions, but a dual circle becomes arbitrarily large as the original
circle becomes arbitrarily small.  In that limit it is effectively
noncompact.  The resulting branes can therefore extend along or be
transverse to directions that appear noncompact at accessible scales.  The
dimension \(p\) counts all spatial directions tangent to the brane, not only
the compact ones used in one derivation.}
\end{classroomqa}

A D0-brane is a new kind of particle.  An open string may run from one
D0-brane to another, return to the same D0-brane, or stretch toward a very
distant brane.  A string taken literally to spatial infinity has infinite
length and hence infinite energy, so the physically useful version places
its other endpoint on another object.

The next case, a D1-brane, is itself string-shaped, but it is not the
fundamental string with which we began.  It can bend and close on itself,
and fundamental strings may end on it.  Its tension scales differently with
the string coupling:
\begin{equation}
  T_{\rm D1}=\frac{1}{2\pi\alpha' g_s},
  \qquad
  T_{\rm F1}=\frac{1}{2\pi\alpha'}.
  \label{eq:l10-dstring-tension}
\end{equation}
At weak coupling \(g_s\ll1\), the D-string is much heavier than the
fundamental string, which explains why it first appears as a massive anchor.

\begin{figure}[H]
  \centering
  \includegraphics[width=0.76\textwidth]{lecture_10_figure_15.png}
  \caption{The board's D1-brane: a bendable one-dimensional object, distinct
  from the fundamental strings that may end on it.}
  \lectureframe{01:24:20}
  \label{fig:l10-d1-brane}
\end{figure}

\begin{classroomqa}
\audiencequestion{Does T-duality act on the D-branes as well as on the
strings? \lecturetimestamp{01:25:35}}
\lecturerresponse{Yes.  A T-duality parallel to a D\(p\)-brane lowers its
dimension to D\((p-1)\), while a T-duality transverse to it raises the
dimension to D\((p+1)\).  Branes are part of the duality dictionary, not
external supports left unchanged by it.}
\end{classroomqa}

\begin{classroomqa}
\audiencequestion{May directions tangent to a brane still be compact?
\lecturetimestamp{01:27:25}}
\lecturerresponse{Yes.  ``Tangent'' and ``compact'' answer different
questions.  A direction may lie along the brane and still close into a
circle; it may also be noncompact.  What matters for the endpoint is whether
the boundary condition in that direction is Neumann or Dirichlet.}
\end{classroomqa}

A D2-brane is a membrane-like sheet.  A D3-brane has three spatial
dimensions.  In a genuinely higher-dimensional spacetime a D3-brane is a
three-dimensional subspace rather than all of space, even though observers
confined to it may experience it as their entire spatial world.

\begin{figure}[H]
  \centering
  \includegraphics[width=0.76\textwidth]{lecture_10_figure_16.png}
  \caption{A D2-brane drawn as a sheet on which string endpoints can move.}
  \lectureframe{01:27:55}
  \label{fig:l10-d2-brane}
\end{figure}

The derivation given here is intentionally a sketch.  It is not by itself a
proof that string theory describes nature.  Within string theory, however,
the consistency checks are numerous: dualities, quantum calculations, and
mathematical results all point to D-branes as genuine dynamical objects.

\section{Open Strings Become Fields on a Brane}

Take an otherwise empty D2-brane and place short open strings on it.  Each
endpoint is confined to the sheet, so the whole low-energy excitation moves
along the sheet.  Two oriented strings whose adjacent endpoints carry
compatible arrows can join.  The joined string may later split again.  Read
as particle histories, these are the vertices of Feynman diagrams: two
particles enter, an intermediate excitation propagates, and two particles
emerge.

\begin{figure}[H]
  \centering
  \includegraphics[width=0.76\textwidth]{lecture_10_figure_17.png}
  \caption{Oriented open strings on a D2-brane joining through compatible
  endpoints.  At low energy the process is read as a field-theory
  interaction on the brane.}
  \lectureframe{01:29:55}
  \label{fig:l10-open-strings-on-brane}
\end{figure}

At energies well below the string scale, the internal vibrational tower is
not excited.  The remaining massless open-string state is a vector field on
the brane worldvolume.  For one brane this is an abelian \(U(1)\) gauge
field, the brane analogue of electromagnetism.  The field does not propagate
through the full ten-dimensional bulk; its quanta are open strings whose
endpoints keep them on the brane.

\section{A Stack of Branes Produces Nonabelian Color}

Now place several parallel branes near one another.  For three branes, label
them Red, Green, and Blue.  An oriented string has an ordered pair of
endpoint labels: it may run from Red to Green, Green to Blue, Red to Red,
and so on.  With \(N\) branes the massless open-string states can be arranged
as an \(N\times N\) matrix
\begin{equation}
  A_\mu{}^{i}{}_{j},
  \qquad i,j=1,\ldots,N.
  \label{eq:l10-chan-paton-matrix}
\end{equation}
These are the Chan--Paton labels of a \(U(N)\) gauge theory.

\begin{figure}[H]
  \centering
  \includegraphics[width=0.76\textwidth]{lecture_10_figure_18.png}
  \medskip

  \begin{adjustbox}{max width=\linewidth}
  \begin{tikzpicture}[x=1cm,y=1cm,>=Latex]
    \draw[very thick,red!70!black] (-4.8,2.0)--(4.8,2.0);
    \draw[very thick,green!55!black] (-4.8,0.8)--(4.8,0.8);
    \draw[very thick,blue!70!black] (-4.8,-0.4)--(4.8,-0.4);
    \node[left] at (-4.8,2.0) {Red};
    \node[left] at (-4.8,0.8) {Green};
    \node[left] at (-4.8,-0.4) {Blue};
    \draw[->,thick] (-3.4,2.0) .. controls (-2.8,1.35) and (-2.3,1.15)
      .. (-1.7,0.8);
    \node at (-2.7,1.05) {$R\!G$};
    \draw[->,thick] (-0.6,0.8) .. controls (0.0,0.25) and (0.6,-0.05)
      .. (1.2,-0.4);
    \node at (0.1,0.0) {$G\!B$};
    \draw[->,thick] (2.1,2.0) .. controls (2.9,2.55) and (3.8,2.55)
      .. (4.3,2.0);
    \node at (3.2,2.65) {$R\!R$};
  \end{tikzpicture}
  \end{adjustbox}
  \caption{Three parallel branes on the board and a clean reconstruction of
  the ordered endpoint labels carried by strings between them.}
  \lectureframe{01:33:15}
  \label{fig:l10-three-brane-stack}
\end{figure}

\begin{classroomqa}
\audiencequestion{Can a string really run from the Red brane to the Green
brane without being blocked by the other drawings? \lecturetimestamp{01:35:07}}
\lecturerresponse{Yes.  The branes are embedded in a higher-dimensional
space.  The two-dimensional board picture is only a projection, so a string
may pass around another projected line in a transverse direction.}
\end{classroomqa}

\begin{classroomqa}
\audiencequestion{Must a Red--Green string combine with a Green--Blue
string in order to become Red--Blue? \lecturetimestamp{01:35:39}}
\lecturerresponse{That is the basic local joining rule.  Compatible
endpoints meet, and the two segments become one string whose unjoined
endpoints retain their Red and Blue labels.  More elaborate processes are
built by repeating the same operation.}
\end{classroomqa}

The joining rule is matrix multiplication made geometric.  A Red--Green
string can join a Green--Blue string because the Green endpoint of one is
compatible with the Green endpoint of the other.  The result is a
Red--Blue string:
\begin{equation}
  (R,G)\circ(G,B)\longrightarrow(R,B).
  \label{eq:l10-color-composition}
\end{equation}
The repeated index is joined and disappears internally.  This is the same
index structure that appears at a Yang--Mills interaction vertex.

\begin{figure}[H]
  \centering
  \includegraphics[width=0.76\textwidth]{lecture_10_figure_19.png}
  \caption{Joining open strings on a stack.  Endpoint orientation enforces
  the same ordered-index composition used by nonabelian gauge fields.}
  \lectureframe{01:36:30}
  \label{fig:l10-color-joining-board}
\end{figure}

For three coincident branes the complete open-string gauge group is
\(U(3)\), with nine vector states.  Its algebra decomposes as
\begin{equation}
  \mathfrak u(3)=\mathfrak{su}(3)\oplus\mathfrak u(1).
  \label{eq:l10-u3-decomposition}
\end{equation}
The eight traceless combinations have the color structure of the QCD
gluons, while the diagonal center-of-mass \(U(1)\) is a separate abelian
field.

\editorialnote{The lecture describes one of the nine color-pair
combinations as able to detach and disappear, thereby leaving eight gluons.
The precise D-brane statement is that a stack of three coincident branes
gives \(U(3)\), not only \(SU(3)\).  Obtaining the QCD octet requires
separating, projecting out, or otherwise decoupling the diagonal
\(U(1)\).  The lecture's picture correctly motivates the ordered color
indices, but it is not a complete derivation of QCD.}

Quark-like matter carries only one color index.  In the brane picture it is
represented by a string with one endpoint on the color stack and the other
on a different, distant brane.  Relative to the color stack, only one end
is visible as a color label.  Reversing the orientation gives the
antiparticle.  A quark-like and antiquark-like string with compatible
endpoints can join and leave the color stack; two identically oriented
quark-like strings cannot annihilate in that way.

The resulting low-energy theories are generally supersymmetric Yang--Mills
theories, with additional matter determined by the branes and their
intersections.  They are close enough to QCD to illuminate nonabelian
dynamics, but they are not automatically the Standard Model.  With only one
brane, the same construction reduces to a \(U(1)\) theory: the short
brane--brane strings are photon-like, while oriented strings reaching a
second sector play the roles of positive and negative electric charges.

\section{D-Strings End as Magnetic Monopoles}

There are now two kinds of string-like object in the story.  The original
fundamental string, or F-string, carries electric flux on a D3-brane.  The
D1-brane, or D-string, is much heavier at weak coupling.  A D-string can also
end on a D3-brane.  Its endpoint is not electrically charged under the
worldvolume gauge field; it is magnetically charged.  From the perspective
of observers on the D3-brane, the endpoint is a magnetic monopole.

\begin{figure}[H]
  \centering
  \includegraphics[width=0.76\textwidth]{lecture_10_figure_20.png}
  \medskip

  \begin{adjustbox}{max width=\linewidth}
  \begin{tikzpicture}[x=1cm,y=1cm,>=Latex]
    \draw[fill=blue!7] (-3.8,-0.5)--(2.6,-0.5)--(3.6,0.35)--(-2.8,0.35)--cycle;
    \node at (2.25,-0.95) {D3-brane};
    \draw[very thick,red!70!black] (0.0,3.0)--(0.0,0.0);
    \fill[black] (0.0,0.0) circle (0.08);
    \node[right] at (0.15,2.1) {D-string};
    \foreach \a in {25,65,115,155} {
      \draw[->,blue!65!black] (0.0,0.0)--({1.25*cos(\a)},{0.65*sin(\a)});
    }
    \node[align=center] at (4.8,1.7) {magnetic flux\\on the D3 worldvolume};
    \draw[->] (3.7,1.35)--(1.05,0.45);
  \end{tikzpicture}
  \end{adjustbox}
  \caption{The board endpoint construction and its clean interpretation.  A
  D-string ending on a D3-brane appears as a magnetic monopole to the
  worldvolume gauge theory.}
  \lectureframe{01:41:30}
  \label{fig:l10-dstring-monopole}
\end{figure}

The mass hierarchy also has the right qualitative form.  Since the D-string
tension carries a factor \(1/g_s\), its endpoint is much heavier than an
ordinary electric excitation at weak coupling.  The electric/magnetic
relation between F-string and D-string endpoints is part of the wider
duality structure of type IIB string theory.

\begin{classroomqa}
\audiencequestion{Is the endpoint of the D-string neutral?
\lecturetimestamp{01:42:19}}
\lecturerresponse{It is neutral with respect to the electric charge carried
by a fundamental-string endpoint, but it is not uncharged.  It carries
magnetic charge and is therefore a monopole in the D3-brane gauge theory.}
\end{classroomqa}

\section{What the Construction Does and Does Not Claim}

The path from T-duality to branes was assembled over decades, not in a
single calculation.  A theory first proposed to describe hadrons became a
candidate theory of quantum gravity.  Its dualities then forced D-branes
into the spectrum, and gauge theories reappeared as the low-energy dynamics
of strings ending on those branes.  In that sense the subject made a full
circle.

D-brane and related dual constructions have become powerful tools for
studying supersymmetric gauge theories, strongly coupled systems, and
connections between gravity and field theory.  They have also informed
models of fluids and plasmas through gauge/gravity duality.  This success is
not the same as deriving the observed particle spectrum directly.  The
compactification geometry, brane configuration, supersymmetry breaking, and
many other ingredients still have to be selected.

\begin{classroomqa}
\audiencequestion{What determines which compact dimension is dualized?
\lecturetimestamp{01:44:59}}
\lecturerresponse{Mathematically, one may dualize any compact direction for
which the required symmetry is present, choosing the description that is
most convenient.  Physically, the brane orientations and compactification
data realized in a particular state are selected by the history and
dynamics of that state.}
\end{classroomqa}

\begin{classroomqa}
\audiencequestion{Which brane supports the monopole construction?
\lecturetimestamp{01:45:21}}
\lecturerresponse{The world experienced by the gauge theory is a D3-brane.
The D-string extends into a transverse direction and ends on that D3-brane;
its endpoint is the monopole seen in the three spatial dimensions of the
brane.}
\end{classroomqa}

\begin{classroomqa}
\audiencequestion{Are D-strings oriented as well?
\lecturetimestamp{01:45:55}}
\lecturerresponse{Yes.  Reversing a D-string's orientation reverses the
corresponding charge.  The two endpoint orientations distinguish monopole
from antimonopole just as orientation distinguishes electric charge for the
fundamental-string endpoints.}
\end{classroomqa}

\begin{classroomqa}
\audiencequestion{Can this simple toroidal picture account for the
left--right asymmetry of chiral particle physics?
\lecturetimestamp{01:46:05}}
\lecturerresponse{Not by itself.  A flat torus is too symmetric.  Chiral
spectra require more structure, such as asymmetric brane intersections,
orbifold projections, or compactification on suitable Calabi--Yau
geometries.  String theory has mechanisms for chirality, but they lie beyond
this elementary T-duality construction.}
\end{classroomqa}

We have followed one symmetry through four concrete forms:
an exchange of spectra, an exchange of worldsheet derivatives, an exchange
of spacetime fields, and finally an exchange of endpoint boundary
conditions.  The last form is the decisive one.  It turns a mathematical
equivalence into the discovery of D-branes, and it turns the dynamics of
open strings into gauge theory on their worldvolumes.  At that point, in
Susskind's phrase, the lecture has used up its time, energy, and momentum.
